% TOP_PROBLEM: {"id": 30005247, "title": "Knots–Quivers Correspondence for All Knots", "category_id": 16, "category": {"id": 16, "name": "physics", "display_name": "Mathematical Physics", "description": "Problems at the intersection of mathematics and physics.", "slug": "physics", "order_index": 16, "created_at": "2026-07-31T15:26:25.671Z"}, "status": "partially_solved", "statusQualification": "The substitutions and framing conventions are part of the correspondence. Equality of one uncolored polynomial would not establish the stated all-color generating-series claim."}
% FIRST_POSTED: 2026-10-11
% ENTRY_KIND: statement_only
% SUBMISSION: {"kind": "statement_only", "category": "top_problems", "problem_id": "30005247", "proposal_type": "established", "source": "Mini-Workshop: Quantization of Complex Symplectic Varieties (2022), Oberwolfach Reports, DOI 10.4171/owr/2022/45, pp. 2629--2658. Original report and mathematical context. https://doi.org/10.4171/owr/2022/45", "status_evidence": "Source status retained from the supplied canonical registry. Dated literature assessment (2026-08-22): The knots--quivers identity is proved for many families and examples, but current literature still describes the all-knots assertion as conjectural, especially beyond the exponential-growth setting. Primary report inspected for formulation; no comprehensive current-status literature review is claimed.", "status_checked_on": "2026-10-11", "duplicate_check": "Canonical numeric ID and public code checked against full registry; absent from baseline selected website registry. Source scope compared with adjacent source questions and existing related catalogue entries. No new identity allocated.", "attribution": "Original authors and contributors credited in Mini-Workshop: Quantization of Complex Symplectic Varieties (2023), Oberwolfach Reports, DOI 10.4171/owr/2022/45, pp. 2629--2658; UnsolvedMath curation credited separately. Definitions and English exposition prepared with Codex assistance; no independent human mathematical verification claimed."}
% SOURCE_STATEMENT: "Does the knots–quivers correspondence hold for every knot, relating its colored HOMFLY–PT generating series $P(x,a,q)$ to the motivic generating series $P_C(x,q)$ of a corresponding quiver?"
% Preparation and final self-audit: CONTRIBUTING.md prompts C and D; source/audit artifacts retained in the local physics_statements_2026_10_11 batch.
% Imported from: physics_statements.tex; UnsolvedMath ID 30005247
% !TeX program = lualatex
\documentclass[11pt]{article}
\usepackage[margin=25mm]{geometry}
\usepackage{fontspec}
\setmainfont{Latin Modern Roman}
\usepackage{amsmath,amssymb,mathtools}
\usepackage{unicode-math}
\setmathfont{Latin Modern Math}
\usepackage{xurl,hyperref}
\hypersetup{colorlinks=true,urlcolor=blue}
\setcounter{secnumdepth}{0}
\setlength{\parindent}{0pt}
\setlength{\parskip}{5pt}
\setlength{\emergencystretch}{3em}
\newcommand{\sref}[2]{\hyperlink{source:#1:#2}{[S#2]}}
\newcommand{\cataloguescope}[1]{\paragraph{Recorded target.}#1}
\newcommand{\R}{\mathbb R}
\newcommand{\C}{\mathbb C}
\newcommand{\N}{\mathbb N}
\newcommand{\Z}{\mathbb Z}
\newcommand{\Q}{\mathbb Q}
\newcommand{\F}{\mathbb F}
\newcommand{\eps}{\varepsilon}
\begin{document}
% UnsolvedMath ID 30005247; source catalogue code OWR-11695854-003
\setcounter{section}{30005246}
\section{Knots–Quivers Correspondence for All Knots}\label{30005247}
{\small UnsolvedMath ID 30005247 · Mathematical Physics}
\subsection{Definitions and mathematical statement}

\paragraph{Shared notation.}$\mathbb N=\{1,2,\ldots\}$; the standard number fields have their usual meanings. All additional conventions are stated locally.


For a knot $K$, colored HOMFLY--PT invariants are assembled into a generating series $P_K(x,a,q)$. A quiver is a finite directed graph, encoded by a matrix $C$; its motivic generating function $P_C$ packages its representation-theoretic invariants. The knots--quivers correspondence identifies these series after the source's grading substitutions.

\paragraph{Catalogue wording retained for traceability.}
Does the knots–quivers correspondence hold for every knot, relating its colored HOMFLY–PT generating series $P(x,a,q)$ to the motivic generating series $P_C(x,q)$ of a corresponding quiver?

\paragraph{Scope and source qualification.}The substitutions and framing conventions are part of the correspondence. Equality of one uncolored polynomial would not establish the stated all-color generating-series claim.
\subsection{Short English statement}
Determine whether every knot has a quiver whose motivic series reproduces all its colored HOMFLY--PT invariants.
\subsection{Sources}
\begin{itemize}
\item[\hypertarget{source:30005247:1}{S1}] ULAM.AI and the UnsolvedMath Contributors, UnsolvedMath, record 30005247 (OWR-11695854-003). Catalogue curation under CC BY 4.0; definitions and English exposition added locally.
\url{https://huggingface.co/datasets/ulamai/UnsolvedMath}.
\item[\hypertarget{source:30005247:2}{S2}] Mini-Workshop: Quantization of Complex Symplectic Varieties (2022), Oberwolfach Reports, DOI 10.4171/owr/2022/45, pp. 2629--2658. Original report and mathematical context.
\url{https://doi.org/10.4171/owr/2022/45}.
\end{itemize}
\label{end:30005247}
\end{document}
